Este trabalho cumpriu seu objetivo principal ao desenvolver, integrar e avaliar uma plataforma robótica física de baixo custo voltada à execução de algoritmos de SLAM visual embarcado. A pesquisa demonstrou, na prática, as viabilidades e os limites de um pipeline robótico executado sobre hardware computacional restrito, como o Raspberry Pi, conjugado a sensores de baixo custo, como a câmera Kinect RGB-D e a IMU MPU-6050.

Diferentemente de pesquisas centradas exclusivamente em simulação, este projeto entregou um protótipo físico funcional capaz de adquirir dados brutos do ambiente, publicar transformações espaciais coerentes no ROS 2, acionar a base móvel e gerar odometria visual em tempo real. Os ensaios experimentais confirmaram que a arquitetura atingiu uma linha de base operacional para o micro-mapeamento, viabilizando a construção de estruturas iniciais de mapas em trajetórias curtas e controladas. 

Os resultados técnicos apontaram que, sob condições ideais de iluminação e sincronização de dados (\texttt{approx\_sync\_max\_interval:=0.025}), o sistema manteve o rastreamento visual e processou nuvens de pontos sem perdas críticas de performance no SBC. Por outro lado, a pesquisa identificou limitações substanciais inerentes ao hardware escolhido. A ausência de odometria de rodas confiável (devido à imprecisão dos encoders originais) exigiu que o pipeline operasse exclusivamente com odometria visual e dados inerciais. Essa configuração revelou-se sensível a movimentos mais agressivos e giros fechados, que ocasionalmente provocavam desfoque nas imagens e subsequente perda de pose (\textit{registration failed}).

Portanto, conclui-se que a execução de SLAM visual em hardware de baixo custo é viável para aplicações de micro-mapeamento e exploração supervisionada, mas ainda impõe desafios severos para alcançar a autonomia completa. O projeto confirmou que gargalos não residem unicamente no processamento (CPU/RAM), mas predominantemente na confiabilidade da percepção visual associada à robustez mecânica do robô.

\section{Trabalhos Futuros}

A plataforma consolidada neste trabalho constitui um artefato valioso para experimentações futuras no contexto da robótica acadêmica. Como desdobramentos diretos desta pesquisa, propõem-se os seguintes trabalhos futuros:

\begin{itemize}
    \item \textbf{Manutenção e Integração da Odometria de Rodas:} Consolidar a leitura dos encoders físicos e implementar filtros de fusão sensorial (como o \textit{robot\_localization}) combinando a odometria visual, dados inerciais e cinemática de rodas, visando mitigar a perda de rastreamento em giros e trajetórias longas;
    \item \textbf{Integração da Pilha de Navegação (Nav2):} Uma vez garantida uma odometria robusta, o passo seguinte consiste em integrar o algoritmo de planejamento global e local para que o robô seja capaz de navegar autonomamente pelos mapas 2D construídos;
    \item \textbf{Descarregamento de Processamento (\textit{Offloading}):} Investigar na prática o ganho de desempenho caso a etapa pesada do SLAM (RTAB-Map) seja transmitida via rede para um servidor externo, deixando o Raspberry Pi responsável apenas pela percepção e navegação de baixo nível;
    \item \textbf{Avaliação com Sensores Modernos:} Substituir o sensor RGB-D legado (Kinect v1) por câmeras estéreo ou LiDARs de baixo custo para avaliar comparativamente a qualidade do mapa e o alívio computacional.
\end{itemize}

O Talus-Droid atendeu à proposta metodológica de \textit{Design Science Research}, fornecendo conhecimentos práticos e uma plataforma testável que contribuem diretamente para a democratização da pesquisa em robótica móvel e navegação autônoma em ambientes fechados.
